\documentclass{article}
\usepackage{graphicx} % Required for inserting images
\usepackage{amssymb}
\usepackage{bbm}
\usepackage{dsfont}

\newcommand{\properideal}{%
\mathrel{\ooalign{$\lneq$\cr\raise.22ex\hbox{$\lhd$}\cr}}}

\usepackage[margin=1in]{geometry} 
\usepackage{amsmath,amsthm,amssymb}
 
\newcommand{\N}{\mathbb{N}}
\newcommand{\Z}{\mathbb{Z}}
 
\newenvironment{theorem}[2][Theorem]{\begin{trivlist}
\item[\hskip \labelsep {\bfseries #1}\hskip \labelsep {\bfseries #2.}]}{\end{trivlist}}
\newenvironment{lemma}[2][Lemma]{\begin{trivlist}
\item[\hskip \labelsep {\bfseries #1}\hskip \labelsep {\bfseries #2.}]}{\end{trivlist}}
\newenvironment{exercise}[2][Exercise]{\begin{trivlist}
\item[\hskip \labelsep {\bfseries #1}\hskip \labelsep {\bfseries #2.}]}{\end{trivlist}}
\newenvironment{problem}[2][Problem]{\begin{trivlist}
\item[\hskip \labelsep {\bfseries #1}\hskip \labelsep {\bfseries #2.}]}{\end{trivlist}}
\newenvironment{question}[2][Question]{\begin{trivlist}
\item[\hskip \labelsep {\bfseries #1}\hskip \labelsep {\bfseries #2.}]}{\end{trivlist}}
\newenvironment{corollary}[2][Corollary]{\begin{trivlist}
\item[\hskip \labelsep {\bfseries #1}\hskip \labelsep {\bfseries #2.}]}{\end{trivlist}}


\begin{document}
\noindent \textbf{MATH 541 Midterm 3}: \\\\
\noindent \textbf{Problem 1} \\
To show that every element $n \in N$ acts trivially on $G / H$, we need to show that for any $n \in N$ and any $x H \in G / H$, we have $n \cdot(x H) = x H$.\\\\
Let $n \in N$ and $x H \in G / H$. We want to show that $n \cdot(x H) = x H$. The action of $G$ on $G / H$ is defined as $g \cdot(x H) = (gx) H$. So, we have:
$$
n \cdot(x H)=(n x) H
$$
Since $N$ is a normal subgroup of $G$ and $N \leq H$, by definition of a normal subgroup, for all $n \in N$ and $g \in G$, we have $g n g^{-1} \in N$. In particular, for the element $x \in G$, we have $x n x^{-1} \in N$. Then, multiplying both sides by $x$ on the right, we get:
$$
x n x^{-1} x=x n \in N
$$
Now, since $N \leq H$, we know that $x n \in H$. As a result, we have $(n x) H = x H$ because the element $x n$ is in $H$. Therefore, we have shown that $n \cdot(x H) = x H$ for any $n \in N$ and any $x H \in G / H$, which means that every element $n \in N$ acts trivially on $G / H$.\\\\
\\
\noindent \textbf{Problem 2} \\




\noindent \textbf{Problem 3} \\
(a) We want to show that $I$ is an ideal of $\mathbb{R}[x]$. We need to verify the following properties:\\\\
1. $I$ is non-empty: Since $(x-1)^2 \in \mathbb{R}[x]$ satisfies $(x-1)^2(1) = (x-1)^2(1)' = 0$, we have $(x-1)^2 \in I$ and thus $I$ is non-empty.\\\\
2. Closed under addition: Let $f(x), g(x) \in I$. Then $f(1) = f'(1) = g(1) = g'(1) = 0$. We want to show that $(f+g)(1) = (f+g)'(1) = 0$. Since $(f+g)(x) = f(x) + g(x)$, we have $(f+g)(1) = f(1) + g(1) = 0 + 0 = 0$. Also, $(f+g)'(x) = f'(x) + g'(x)$, so $(f+g)'(1) = f'(1) + g'(1) = 0 + 0 = 0$. Thus, $f(x) + g(x) \in I$.\\\\
3. Closed under multiplication by elements from $\mathbb{R}[x]$: Let $f(x) \in I$ and $g(x) \in \mathbb{R}[x]$. We want to show that $(fg)(1) = (fg)'(1) = 0$. We have $(fg)(1) = f(1)g(1) = 0 \cdot g(1) = 0$. To find the derivative, we use the product rule: $(fg)'(x) = f'(x)g(x) + f(x)g'(x)$. Then, $(fg)'(1) = f'(1)g(1) + f(1)g'(1) = 0 \cdot g(1) + 0 \cdot g'(1) = 0$. Thus, $f(x)g(x) \in I$.\\\\
Since $I$ is non-empty and closed under addition and multiplication by elements from $\mathbb{R}[x]$, $I$ is an ideal of $\mathbb{R}[x]$.\\\\

\noindent(b) To show that $I$ is generated by $(x-1)^2$, we will demonstrate two things:\\\\
1. $(x-1)^2 \in I$.\\
2. Every element of $I$ can be written as a multiple of $(x-1)^2$.\\\\
First, let's show that $(x-1)^2 \in I$. Observe that $(x-1)^2$ has the property that $f(1)=0$ and $f'(1)=0$. Thus, $(x-1)^2 \in I$.\\\\
Now, let $f(x) \in I$. This means that $f(1) = 0$ and $f'(1) = 0$. We can write $f(x) = (x-1)g(x) + r(x)$, where $g(x), r(x) \in \mathbb{R}[x]$ and $r(x) = 0$ or $\deg(r) < \deg(x-1) = 1$. Since $\mathbb{R}[x]$ is a Euclidean domain, such a division is always possible.\\\\
Now, let's analyze $r(x)$. Since $r(x)$ has degree less than 1, it can be written as $r(x) = a_0 + a_1(x-1)$, where $a_0, a_1 \in \mathbb{R}$. By the properties of $I$, we know that $f(1) = 0$, so
$$
0=f(1)=(1-1) g(1)+a_0+a_1(1-1)=a_0
$$
Thus, $a_0 = 0$ and $r(x) = a_1(x-1)$.\\\\
Now, let's take the derivative of $f(x)$:
$$
f^{\prime}(x)=(x-1) g^{\prime}(x)+g(x)+a_1
$$
Since $f'(1) = 0$, we have
$$
0=(1-1) g^{\prime}(1)+g(1)+a_1=g(1)+a_1
$$
which implies that $a_1 = -g(1)$.\\\\
Now we can rewrite $r(x)$ as $r(x) = -g(1)(x-1)$. Therefore, $f(x) = (x-1)(g(x) - g(1))$. Notice that this means $f(x)$ is a multiple of $(x-1)^2$, as $(x-1)(g(x) - g(1))$ can be written as $(x-1)^2 h(x)$, where $h(x) = \frac{g(x) - g(1)}{x-1}$.\\\\
We have shown that every element of $I$ can be written as a multiple of $(x-1)^2$. Thus, $I$ is generated by $(x-1)^2$.\\\\



\noindent \textbf{Problem 4} \\
To prove this, we will use the Chinese Remainder Theorem for the ring $\mathbb{Z}[\sqrt{5}]$. We need to show that $(8, \sqrt{5}) = (1)$, i.e., the ideals generated by $8$ and $\sqrt{5}$ have no common elements other than $0$.\\\\
Suppose there exists a non-zero element $\delta \in \mathbb{Z}[\sqrt{5}]$ such that $\delta \in (8, \sqrt{5})$. Then, there exist elements $x, y \in \mathbb{Z}[\sqrt{5}]$ such that $\delta = 8x = \sqrt{5}y$. We can write $x = a + b\sqrt{5}$ and $y = c + d\sqrt{5}$ for some integers $a, b, c, d \in \mathbb{Z}$. Then, we have $\delta = 8(a + b\sqrt{5}) = \sqrt{5}(c + d\sqrt{5})$.\\\\
Comparing the coefficients of $\sqrt{5}$ on both sides, we get $8b = c$ and $8a = 5d$. Since $8b = c$, $c$ is a multiple of 8. Since $8a = 5d$, 5 divides $8a$, and since 5 is prime and does not divide 8, it must divide $a$. Let $a = 5a'$ and $c = 8c'$. Then $d = 8a'$. Now, we can write $y = c + d\sqrt{5} = 8c' + 8a'\sqrt{5} = 8(c' + a'\sqrt{5})$.\\\\
So, we have $8x = 8(a + b\sqrt{5}) = \delta = \sqrt{5}(8(c' + a'\sqrt{5}))$. Canceling the 8's, we get $x = \sqrt{5}(c' + a'\sqrt{5})$. This implies that $\sqrt{5}$ divides $x$, contradicting the fact that $\delta$ is a common divisor of $8x$ and $\sqrt{5}y$ and the assumption that $\delta$ is non-zero. Thus, $(8, \sqrt{5}) = (1)$.\\\\
Now, by the Chinese Remainder Theorem, there exists an element $\gamma \in \mathbb{Z}[\sqrt{5}]$ such that $\gamma \equiv \alpha \bmod 8$ and $\gamma \equiv \beta \bmod \sqrt{5}$. In other words, $8 \mid \gamma - \alpha$ and $\sqrt{5} \mid \gamma - \beta$.\\\\


\noindent \textbf{Problem 5 (draft)} \\
To show that $(2,1+\sqrt{3})$ is a maximal ideal in $\mathbb{Z}[\sqrt{3}]$, we can use the hint and follow a similar approach to what we did for $\mathbb{Z}[\sqrt{-5}]$. We will consider the quotient ring $\mathbb{Z}[\sqrt{3}]/(2,1+\sqrt{3})$ and show that it is a field. If it is a field, then $(2,1+\sqrt{3})$ is a maximal ideal according to the definition.\\\\
Let's define a homomorphism $\phi: \mathbb{Z}[\sqrt{3}] \to \mathbb{Z}_2[x]/(x^2-3)$ as follows:
$$
\phi(a+b \sqrt{3})=\bar{a}+\bar{b} x(\bmod 2)
$$
where $a,b \in \mathbb{Z}$, and $\overline{a}$ and $\overline{b}$ are the residues of $a$ and $b$ modulo 2, respectively. It's easy to see that this map is a homomorphism.\\\\
Now, notice that $\phi(2)=\overline{2}=\overline{0}$, and $\phi(1+\sqrt{3})=\overline{1}+\overline{1}x=x+1$. Thus, $\phi$ maps the ideal $(2,1+\sqrt{3})$ to the ideal $(0, x+1)$ in $\mathbb{Z}_2[x]/(x^2-3)$. Then, the kernel of $\phi$ is $(2,1+\sqrt{3})$, and by the First Isomorphism Theorem, we have:
$$
\mathbb{Z}[\sqrt{3}] /(2,1+\sqrt{3}) \cong \phi(\mathbb{Z}[\sqrt{3}]) \subseteq \mathbb{Z}_2[x] /\left(x^2-3\right)
$$
Now, we want to show that $\mathbb{Z}_2[x]/(x^2-3)$ is a field. In $\mathbb{Z}_2$, we have that $-1 = 1$. Therefore, in $\mathbb{Z}_2[x]$, the polynomial $x^2-3$ can be rewritten as $x^2+1$. Now, we can see that $x^2+1$ is an irreducible polynomial in $\mathbb{Z}_2[x]$, since it has no roots in $\mathbb{Z}_2$ (neither $0$ nor $1$ are roots). Thus, the ideal $(x^2+1)$ is a maximal ideal in $\mathbb{Z}_2[x]$, which implies that the quotient ring $\mathbb{Z}_2[x]/(x^2+1)$ is a field.\\\\
Since $\mathbb{Z}[\sqrt{3}]/(2,1+\sqrt{3})$ is isomorphic to a subring of $\mathbb{Z}_2[x]/(x^2-3)$, which is a field, we can conclude that $\mathbb{Z}[\sqrt{3}]/(2,1+\sqrt{3})$ is also a field. Therefore, $(2,1+\sqrt{3})$ is a maximal ideal in $\mathbb{Z}[\sqrt{3}]$.\\\\



\noindent \textbf{Problem 6} \\
Consider the ring morphism $\varphi: \mathbb{Z} \rightarrow \mathbb{Z}[x]$ defined by $\varphi(n) = n$ for all $n \in \mathbb{Z}$. This is a valid ring morphism as it preserves addition, multiplication, and the identity elements.\\\\
Now, let $\mathfrak{m} = (x) \subseteq \mathbb{Z}[x]$, which is a maximal ideal in $\mathbb{Z}[x]$.\\\\
Let's compute the pre-image of $\mathfrak{m}$ under $\varphi$:
$$\varphi^{-1}(\mathfrak{m}) = {n \in \mathbb{Z} \mid \varphi(n) \in (x)}$$
Since $\varphi(n) = n$ and the only constant polynomial in $(x)$ is the zero polynomial, we have $\varphi^{-1}(\mathfrak{m}) = (0)$, which is the ideal consisting of only the zero element in $\mathbb{Z}$.\\\\
The pre-image $\varphi^{-1}(\mathfrak{m})$ is not a maximal ideal in $\mathbb{Z}$, because $\mathbb{Z}/(0) \cong \mathbb{Z}$ is not a field. Moreover, it is not equal to the entire ring $R=\mathbb{Z}$, as it contains only the zero element.\\\\
\end{document}